%%%%%%%%%%%%%%%%%%%%%%%%%%%%%%%%%%%%%%%%%%%%%%%%%%%%%%%%%%%%%%%%%%%%%%%%%%%
% Handout und Uebungsblatt zu Lehreinheit 5
% Im Skript am Kapitelende, in handouts.tex als Druckvorlage.
%%%%%%%%%%%%%%%%%%%%%%%%%%%%%%%%%%%%%%%%%%%%%%%%%%%%%%%%%%%%%%%%%%%%%%%%%%%
\begin{handout}{Hypothesen und Falsifizierbarkeit}

\textbf{Worum es heute ging.} Ein Schaltbrett mit vier Schaltern und vier
Lampen: Sie haben ein Modell gebildet und es oft bestätigt gesehen. Dann
genügte ein Umstecken, und das Modell war widerlegt. Viele Bestätigungen
machen eine allgemeine Aussage nicht gewiss; eine einzige Gegenbeobachtung
kann sie zu Fall bringen. Die Lösung des Bretts liegt im Gerät, und es ist
komplizierter, als Ihr Modell erlaubt hat.

\textbf{Kernbegriffe}
\begin{description}
  \item[Induktion / Deduktion / Abduktion] von Fällen auf eine Regel /
    von der Regel auf eine Vorhersage / von einer Überraschung auf die
    beste Erklärung. Logisch zwingend ist nur die Deduktion.
  \item[Falsifizierbarkeit] Eine Aussage ist prüfbar, wenn sich eine
    Beobachtung angeben lässt, die ihr widerspräche
    \parencite{popper1934logik}. Bewährt heißt: Widerlegungsversuche sind
    gescheitert. Bewiesen ist damit nichts.
  \item[Immunisierung] eine Aussage wird so formuliert oder nachträglich
    angepasst, dass nichts sie widerlegen kann.
  \item[H\textsubscript{0} / H\textsubscript{1}] Nullhypothese (kein
    Unterschied, kein Zusammenhang) und Alternativhypothese (Ihre
    Vermutung). Geprüft wird, ob die Daten mit H\textsubscript{0}
    vereinbar sind.
  \item[Graduelle Bestätigung] Tritt eine vorhergesagte Beobachtung $E$
    ein, steigt nur die Wahrscheinlichkeit von $H$, und zwar umso mehr, je
    unwahrscheinlicher $E$ vorher war \parencite{horwich1982probability}.
  \item[Duhem-Quine] Geprüft wird immer die Hypothese \emph{samt}
    Hilfsannahmen; ein Widerspruch sagt nicht, welches Glied falsch ist.
  \item[Congruence Bias] nur die eigene Hypothese direkt testen, nie die
    Alternativen (Wason-2-4-6-Test).
\end{description}

\textbf{Raster: Ist meine Hypothese gut?}

\begin{center}
\small
\begin{tabular}{p{0.2\textwidth}p{0.7\textwidth}}
\toprule
präzise & Sind alle Begriffe operationalisiert? \\
falsifizierbar & Welches Ergebnis würde ihr widersprechen? \\
riskant & Wäre das Ergebnis ohne die Hypothese unwahrscheinlich? \\
begrenzt & Für wen, wann, unter welchen Bedingungen gilt sie? \\
begründet & Woraus folgt sie: Theorie, Literatur, Beobachtung? \\
offen & Sind die Hilfsannahmen benannt? \\
\bottomrule
\end{tabular}
\end{center}

\begin{merkbox}
Fragen Sie nicht \enquote{Was spricht für meine Hypothese?}, sondern
\enquote{Was würde mich umstimmen?} Eine Beobachtung ist nur so viel wert
wie das Risiko, das die Hypothese mit ihr eingegangen ist.
\end{merkbox}

\textbf{Zum Weiterlesen:} \textcite{popper1934logik};
\textcite{wason1960failure}; zur Vertiefung
\textcite{horwich1982probability}.
\end{handout}

\begin{uebungsblatt}{Hypothesen und Falsifizierbarkeit}

\aufgabe{Falsifizierbar oder nicht?}[in der Stunde, Kleingruppe, 10 Min.]
Entscheiden Sie für jede Aussage, ob sie \emph{in der vorliegenden Form}
falsifizierbar ist (ja / nein / erst nach Präzisierung). Falsifizierbar
heißt: Es lässt sich eine Beobachtung angeben, die der Aussage
widerspräche. Geben Sie für jede prüfbare Aussage eine konkrete
Beobachtung an, die sie widerlegen würde.
Formulieren Sie vage Aussagen so um, dass sie prüfbar werden. (Die
Aussagen sind für diese Übung konstruiert.)

\begin{enumerate}
  \item[A1] \enquote{Ehrliche Menschen schlafen besser.}
  \item[A2] \enquote{Studierende, die in der Woche vor der Klausur im
    Mittel mindestens sieben Stunden pro Nacht schlafen, erzielen in
    dieser Klausur im Mittel mehr Punkte als Studierende mit weniger als
    fünf Stunden.}
  \item[A3] \enquote{Alles, was geschieht, geschieht aus einem Grund, auch
    wenn wir ihn nicht immer erkennen können.}
  \item[A4] \enquote{Astrologinnen und Astrologen können einem
    Geburtshoroskop häufiger als zufällig das passende von drei
    Persönlichkeitsprofilen zuordnen.}
  \item[A5] \enquote{Das Mittel wirkt bei allen, die offen dafür sind.
    Wer keine Wirkung spürt, war innerlich nicht bereit.}
\end{enumerate}

\begin{center}
\renewcommand{\arraystretch}{2.4}
\begin{tabular}{|p{0.06\textwidth}|p{0.14\textwidth}|p{0.68\textwidth}|}
\hline
\textbf{Nr.} & \textbf{falsifizierbar?} & \textbf{Widerlegungsbedingung bzw.\ präzisierte Fassung} \\
\hline
A1 & & \\ \hline
A2 & & \\ \hline
A3 & & \\ \hline
A4 & & \\ \hline
A5 & & \\ \hline
\end{tabular}
\end{center}
\end{uebungsblatt}

\begin{uebungsblattb}{2}{Schaltbrett und eigene Hypothese}

\aufgabe{Das Schaltbrett im Rückblick}[in der Stunde oder Hausaufgabe]
\begin{enumerate}
  \item Notieren Sie das Modell, das Ihre Gruppe in Runde~1 gebildet hat,
    und die Zahl der Versuche, die es bestätigt haben.
  \item Welche \emph{eine} Eingabe hätte Ihr Modell am stärksten
    gefährdet? Begründen Sie mit Horwichs Widerlegungsrisiko: Ein Test
    bestätigt umso stärker, je unwahrscheinlicher sein Ausgang ohne das
    Modell wäre.
  \item Nennen Sie zwei Hilfsannahmen, die Sie beim Explorieren
    stillschweigend gemacht haben (Duhem-Quine: Geprüft wird immer das
    Modell \emph{samt} solcher Annahmen).
\end{enumerate}
\schreiblinien{5}

\aufgabe{Null- und Alternativhypothese für Ihre Forschungsfrage}[Beginn in der Stunde, Abgabe bis LE~6]
\begin{center}
\renewcommand{\arraystretch}{1.6}
\begin{tabular}{|p{0.3\textwidth}|p{0.6\textwidth}|}
\hline
Meine Forschungsfrage (aus LE~2) & \\[6mm] \hline
Art der Hypothese (Unterschied / Zusammenhang / Veränderung; gerichtet?) & \\ \hline
H\textsubscript{1} (Alternativhypothese) & \\[8mm] \hline
H\textsubscript{0} (Nullhypothese) & \\[8mm] \hline
Operationalisierung: Begriff $\to$ Messvorschrift & \\[8mm] \hline
Population, Bedingungen, Zeitraum & \\ \hline
Ich verwerfe H\textsubscript{1}, wenn \ldots & \\[8mm] \hline
Hilfsannahmen, die ich dabei mache & \\[6mm] \hline
\end{tabular}
\end{center}

\aufgabe{Tagebuch des Staunens}[Hausaufgabe]
Beschreiben Sie ein Erlebnis dieser Woche, in dem Sie (oder jemand
anderes) eine Annahme nur bestätigt statt geprüft haben. Welcher Versuch
hätte die Annahme wirklich gefährdet?

\begin{lehrhinweis}[Erwartungshorizont]
\textbf{Aufgabe~1.} A1: in dieser Form nicht prüfbar, weil
\enquote{ehrlich} und \enquote{besser schlafen} unbestimmt sind;
präzisierbar, z.\,B.\ \enquote{Erwachsene mit höheren Werten in einem
validierten Fragebogen zur Ehrlichkeit haben über zwei Wochen eine längere,
per Aktivitätstracker gemessene Schlafdauer}; widerlegt, wenn kein oder ein
negativer Zusammenhang auftritt. A2: falsifizierbar; widerlegt, wenn die
Gruppe mit mindestens sieben Stunden im Mittel gleich viele oder weniger
Punkte erzielt. Gute Antworten bemerken, dass ein solcher Befund noch
keine Kausalität zeigt (Brücke zu LE~6). A3: nicht falsifizierbar; der
Zusatz \enquote{auch wenn wir ihn nicht erkennen} schließt jede
Gegenbeobachtung aus. A4: falsifizierbar; widerlegt, wenn die
Trefferquote im Doppelblindversuch nicht über einem Drittel liegt. So
wurde die Aussage tatsächlich geprüft \parencite{carlson1985doubleblind},
mit einem Ergebnis auf Zufallsniveau. A5: immunisiert; jede ausbleibende
Wirkung wird der Person zugeschrieben. Prüfbar wird die Aussage erst, wenn
\enquote{offen dafür} unabhängig und \emph{vor} der Einnahme gemessen wird.
\textbf{Aufgabe~2.} Bewertet wird, ob die gewählte Eingabe einen
Ausgang hat, den das Modell \emph{nicht} sicher vorhersagt, und ob als
Hilfsannahmen z.\,B.\ funktionierende Lampen, volle Batterien, Gleichheit
von sichtbarer und tatsächlicher Verdrahtung genannt werden.
\textbf{Aufgabe~3.} H\textsubscript{0} und H\textsubscript{1} schließen
sich aus; die Verwerfungsbedingung nennt eine Messung und ein Kriterium;
die Operationalisierung deckt alle Begriffe der Hypothese ab.
\end{lehrhinweis}
\end{uebungsblattb}
